% Part: first-order-logic
% Chapter: sequent-calculus
% Section: rules-and-proofs

\documentclass[../../../include/open-logic-section]{subfiles}

\begin{document}

\iftag{FOL}
      {\olfileid{fol}{seq}{rul}}
      {\olfileid{pl}{seq}{rul}}

\olsection{বিধি ও \usetoken{P}{derivation}}

পরের আলোচনায় $\Gamma, \Delta, \Pi, \Lambda$ দিয়ে !!{sentence}s-এর
সসীম অনুক্রম বোঝানো হবে।

\begin{defn}[সিকোয়েন্ট]
কোনো \emph{সিকোয়েন্ট} হলো নিচের রূপের একটি অভিব্যক্তি:
\[
\Gamma \Sequent \Delta
\]
এখানে $\Gamma$ ও $\Delta$ হলো ভাষা~$\Lang L$-এর !!{sentence}s-এর
সসীম (সম্ভবত ফাঁকা) অনুক্রম। $\Gamma$-কে \emph{পূর্বাংশ} এবং
$\Delta$-কে \emph{উত্তরাংশ} বলা হয়।
\end{defn}

\begin{explain}
সিকোয়েন্টের পিছনের স্বজ্ঞাত ধারণাটি হলো: পূর্বাংশের সব !!{sentence}s
সত্য হলে উত্তরাংশের !!{sentence}s-এর মধ্যে অন্তত একটি সত্য। অর্থাৎ, যদি
$\Gamma = \tuple{!A_1, \dots, !A_m}$ এবং
$\Delta = \tuple{!B_1, \dots, !B_n}$ হয়, তবে $\Gamma \Sequent \Delta$
ঠিক তখনই সত্য, যখন
\[
(!A_1 \land \cdots \land !A_m) \lif (!B_1 \lor \cdots \lor
!B_n)
\]
সত্য। দুটি বিশেষ ক্ষেত্র আছে: $\Gamma$ ফাঁকা এবং $\Delta$ ফাঁকা।
$\Gamma$ ফাঁকা, অর্থাৎ $m = 0$ হলে, $\quad
\Sequent \Delta$ ঠিক তখনই সত্য, যখন $!B_1 \lor \dots \lor !B_n$
সত্য। $\Delta$ ফাঁকা, অর্থাৎ $n = 0$ হলে, $\Gamma \Sequent \quad$
ঠিক তখনই সত্য, যখন $\lnot(!A_1 \land \dots \land !A_m)$ সত্য।
সংশ্লিষ্ট !!{sentence}-টি বৈধ হলে এবং কেবল হলেই কোনো সিকোয়েন্টকে
বৈধ বলা হয়।
\end{explain}

$\Gamma$ যদি !!{sentence}s-এর একটি অনুক্রম হয়, তবে~$\Gamma$-র ডান
প্রান্তে $!A$ জুড়ে পাওয়া অনুক্রমকে $\Gamma, !A$ লিখি (আর~$\Gamma$-র বাঁ প্রান্তে
$!A$ জুড়ে পাওয়া অনুক্রমকে $!A, \Gamma$ লিখি)। $\Delta$-ও যদি
!!{sentence}s-এর একটি অনুক্রম হয়, তবে $\Gamma, \Delta$ হলো অনুক্রম
দুটির সংযুক্তি।

\begin{defn}[প্রারম্ভিক সিকোয়েন্ট]
কোনো \emph{প্রারম্ভিক সিকোয়েন্ট} হলো
\iftag{prvFalse,prvTrue}{নিচের কোনো একটি রূপের সিকোয়েন্ট:
  \begin{enumerate}
    \item $!A \Sequent !A$
    \tagitem{prvTrue}{$\quad \Sequent \ltrue$}{}
    \tagitem{prvFalse}{$\lfalse \Sequent \quad$}{}
  \end{enumerate}}
{রূপ~$!A \Sequent !A$-এর সিকোয়েন্ট}; এখানে $!A$ ভাষাটির যেকোনো
!!{sentence}।
\end{defn}

সিকোয়েন্ট কলনে !!^{derivation}s হলো সিকোয়েন্টের বিশেষ ধরনের বৃক্ষ।
তাদের সর্বোচ্চ সিকোয়েন্টগুলি প্রারম্ভিক সিকোয়েন্ট; কোনো সিকোয়েন্ট
অন্য এক বা দুটি সিকোয়েন্টের নিচে থাকলে, সেটিকে কোনো অনুমান-বিধি
সঠিকভাবে প্রয়োগ করে পাওয়া যেতে হবে। $\Log{LK}$-এর বিধিগুলি প্রধানত দুই প্রকার:
\emph{যৌক্তিক} বিধি ও \emph{গঠনগত} বিধি। নিচের সিকোয়েন্টে $!A$ এবং/অথবা
$!B$-যুক্ত !!{sentence}-এর !!{main
  operator} অনুসারে যৌক্তিক বিধিটির
নাম হয়। প্রতিটি বিধির দুটি রূপ থাকে—একটি বাঁদিকে !!{operator}-যুক্ত
!!{sentence} থাকা সিকোয়েন্ট অনুমানের জন্য, অন্যটি ডানদিকে সেই
!!{sentence} থাকা সিকোয়েন্ট অনুমানের জন্য।

\end{document}
